\section{Part 4: Training the Final Nanochat Model}

\subsection{Motivation}

The picochat ablation experiments in Part~2 compared several architectural
modifications, including SwiGLU and LayerScale. Although the baseline
architecture achieved the best results at the pico scale, architectural
changes often behave differently as model capacity increases.

Recent large language models such as PaLM and LLaMA adopt gated MLP
variants (e.g., SwiGLU), suggesting that such modifications may become
beneficial at larger scales. Therefore, in addition to the baseline
configuration, we trained a larger nanochat model incorporating SwiGLU
to examine whether its advantages emerge under increased model capacity.

\subsection{Model Configuration}

The final nanochat models were trained using the configuration shown in
Table~\ref{tab:nano_config}. Compared to picochat, the nanochat model
increases network depth and effective batch size, enabling training on
a substantially larger token budget.

\begin{table}[H]
\centering
\begin{tabular}{lc}
\hline
\textbf{Parameter} & \textbf{Value} \\
\hline
Model Depth & 20 layers \\
Sequence Length & 512 \\
Training Iterations & 4357 \\
Device Batch Size & 32 \\
Total Batch Size & 1,048,576 \\
GPU Type & H100 \\
Precision & FP8 (tensorwise) \\
Window Pattern & L \\
\hline
\end{tabular}
\caption{Nanochat training configuration.}
\label{tab:nano_config}
\end{table}

Training was performed on the full pretraining dataset using
the nanochat training framework.

Two models were trained:

\begin{itemize}
\item \textbf{Nanochat Baseline}: standard ReLU$^2$ MLP.
\item \textbf{Nanochat SwiGLU}: identical architecture with a SwiGLU MLP block.
\end{itemize}

\subsection{Training Procedure}

Training was conducted on H100 GPUs via Modal using FP8 tensorwise
precision to improve throughput and memory efficiency. Training progress
was monitored using Weights \& Biases.

Both models exhibited nearly identical loss trajectories throughout
training, indicating similar optimization behavior. However, the SwiGLU
model showed lower token throughput due to the additional gating
operations within the MLP block.

\subsection{Scaling Prediction}

Before training the nanochat models, we estimate how performance should
change when scaling from picochat (6 layers) to nanochat (20 layers).

The picochat baseline achieved:

\begin{itemize}
\item BPB: 1.0575
\item CORE: 0.0615
\end{itemize}

Increasing model capacity and training compute should significantly
improve language modeling performance. We therefore predict that the
nanochat model will achieve substantially lower BPB and higher CORE
scores than picochat.

After training, the results confirm this prediction. The nanochat
baseline achieves:

\begin{itemize}
\item Validation BPB: 0.7822
\item CORE score: 0.2038
\end{itemize}

This corresponds to roughly a \textbf{26\% reduction in BPB} and more
than a \textbf{threefold increase in CORE}. The improvement is far larger
than the differences between architectural variants, indicating that
model scale is the dominant factor affecting performance.

\subsection{Evaluation}

Following the evaluation protocol in the original nanochat repository,
models are evaluated using two metrics:

\begin{itemize}
\item \textbf{Bits-per-byte (BPB)}, measuring next-token prediction quality.
\item \textbf{CORE benchmark score}, aggregating performance across
22 reasoning and language understanding tasks.
\end{itemize}

\subsubsection{Aggregate Results}

\begin{table}[H]
\centering
\begin{tabular}{lccc}
\hline
\textbf{Model} & \textbf{Train BPB $\downarrow$} & \textbf{Val BPB $\downarrow$} & \textbf{CORE $\uparrow$} \\
\hline
Nanochat Baseline & \textbf{0.782469} & \textbf{0.782195} & 0.2038 \\
Nanochat SwiGLU & 0.786032 & 0.785764 & \textbf{0.2086} \\
\hline
\end{tabular}
\caption{Aggregate nanochat results.}
\label{tab:part4_aggregate_results}
\end{table}

The baseline model achieves the lowest BPB, indicating stronger
token-level language modeling. In contrast, the SwiGLU model achieves a
slightly higher CORE score.

This suggests a small trade-off: SwiGLU does not improve token-level
prediction efficiency but may produce representations that benefit
certain reasoning tasks.

\subsubsection{CORE Benchmark Analysis}

Figure~\ref{fig:part4_core_bar} shows the aggregate CORE scores.

\begin{figure}[H]
\centering
\includegraphics[width=0.55\linewidth]{p4/Screenshot 2026-03-06 113232.png}
\caption{Aggregate CORE scores.}
\label{fig:part4_core_bar}
\end{figure}

\begin{figure}[H]
\centering
\includegraphics[width=\linewidth]{p4/Screenshot 2026-03-06 121848.png}
\caption{Per-task CORE benchmark scores for the baseline and SwiGLU nanochat models.}
\label{fig:part4_core_tasks}
\end{figure}

Task-level results are shown in Table~\ref{tab:part4_task_level_results}.

\begin{table}[H]
\centering
\small
\begin{tabular}{lcc}
\hline
\textbf{Benchmark} & \textbf{Baseline} & \textbf{SwiGLU} \\
\hline
agi\_eval\_lsat\_ar & 0.2348 & \textbf{0.2565} \\
arc\_challenge & \textbf{0.3584} & 0.3456 \\
arc\_easy & 0.6599 & 0.6599 \\
bigbench\_cs\_algorithms & 0.3841 & \textbf{0.4008} \\
bigbench\_dyck\_languages & \textbf{0.1260} & 0.0990 \\
bigbench\_language\_identification & \textbf{0.2520} & 0.2510 \\
bigbench\_operators & \textbf{0.1667} & 0.1619 \\
bigbench\_qa\_wikidata & 0.4816 & \textbf{0.4818} \\
bigbench\_repeat\_copy\_logic & \textbf{0.0312} & 0.0000 \\
boolq & \textbf{0.6043} & 0.5914 \\
commonsense\_qa & 0.2260 & \textbf{0.3522} \\
copa & 0.6200 & \textbf{0.6900} \\
coqa & \textbf{0.2068} & 0.2026 \\
hellaswag & \textbf{0.4409} & 0.4359 \\
hellaswag\_zeroshot & \textbf{0.4573} & 0.4489 \\
jeopardy & 0.0586 & \textbf{0.0709} \\
lambada\_openai & \textbf{0.3868} & 0.3842 \\
openbook\_qa & \textbf{0.3600} & 0.3420 \\
piqa & 0.6763 & \textbf{0.6823} \\
squad & 0.0219 & \textbf{0.0799} \\
winograd & \textbf{0.6667} & 0.6154 \\
winogrande & \textbf{0.5564} & 0.5304 \\
\hline
\end{tabular}
\caption{Task-level CORE benchmark results.}
\label{tab:part4_task_level_results}
\end{table}

The baseline performs slightly better on tasks resembling language
modeling and narrative completion (e.g., HellaSwag, BoolQ). In contrast,
SwiGLU improves several reasoning-oriented tasks such as
\texttt{commonsense\_qa}, \texttt{copa}, and \texttt{squad}. This suggests
that gated MLP layers may influence internal representations used in
reasoning tasks.

\subsubsection{Training Dynamics}

\begin{figure}[H]
\centering

\begin{minipage}{0.9\linewidth}
\centering
\includegraphics[width=\linewidth]{p4/W&B Chart 3_6_2026, 11_33_18 AM.png}
\subcaption{Training Loss}
\end{minipage}

\vspace{0.6em}

\begin{minipage}{0.48\linewidth}
\centering
\includegraphics[width=\linewidth]{p4/W&B Chart 3_6_2026, 11_51_35 AM.png}
\subcaption{Training Time}
\end{minipage}
\hfill
\begin{minipage}{0.48\linewidth}
\centering
\includegraphics[width=\linewidth]{p4/W&B Chart 3_6_2026, 11_51_43 AM.png}
\subcaption{Token Throughput}
\end{minipage}

\caption{Training dynamics comparison between nanochat models.}
\label{fig:training_dynamics}

\end{figure}

Both models exhibit nearly identical loss curves, indicating similar
optimization behavior. However, the baseline maintains higher token
throughput ($\sim160$k tokens/sec) compared to the SwiGLU model
($\sim130$k tokens/sec), resulting in shorter training time.

This difference is also reflected in compute cost. The baseline training
required approximately \$33.80 in Modal compute credits, while the
SwiGLU experiment required \$44.12.

\subsubsection{Comparison with Picochat}

At the pico scale, the baseline architecture outperformed all variants
under both BPB and CORE:

\begin{itemize}
\item Pico baseline: BPB 1.0575, CORE 0.0615
\item Pico SwiGLU: BPB 1.0589, CORE 0.0554
\end{itemize}

At the nano scale, the baseline still achieves better BPB but SwiGLU
slightly surpasses it in CORE. This demonstrates that architectural
effects can change with model scale. While the modification does not
dominate overall, it becomes more competitive as model capacity
increases.

\subsection{Emergent Abilities}

Emergent abilities refer to capabilities that appear in larger models
but are absent or unreliable in smaller ones. To investigate this
phenomenon, we tested both picochat and nanochat models on a set of
diagnostic questions covering factual knowledge, scientific explanation,
and logical reasoning.

We selected questions where nanochat produced meaningful answers while
picochat failed or generated nonsensical output. See Table~\ref{tab:emergent_questions} for the list of questions.

\begin{table}[H]
\centering
\small
\begin{tabular}{p{0.95\linewidth}}
\hline
\textbf{Emergent Ability Questions} \\
\hline
1. What is the capital city of Australia? \\
2. Which country has the city of Barcelona? \\
3. Who wrote the novel \textit{1984}? \\
4. What is the largest ocean on Earth? \\
5. What gas do humans breathe in that plants produce during photosynthesis? \\
6. Why does the sky appear blue during the day? \\
7. Why do objects fall toward the ground on Earth? \\
8. What part of a plant performs photosynthesis? \\
9. What happens to water when it freezes? \\
10. If all cats are animals and all animals breathe air, do cats breathe air? \\
\hline
\end{tabular}
\caption{Questions demonstrating emergent abilities when scaling from
picochat to nanochat.}
\label{tab:emergent_questions}
\end{table}


The picochat model frequently produced incoherent or repetitive
responses. In contrast, both nanochat models correctly answered several
questions, identifying facts such as \textit{Canberra} as Australia's
capital and \textit{oxygen} as the gas produced during photosynthesis.

Between the two nanochat variants, qualitative differences were small.
The baseline model tended to produce longer explanations, whereas the
SwiGLU model often produced shorter and more direct answers.

Overall, these examples illustrate that scaling from picochat to
nanochat enables the model to retrieve factual knowledge and generate
basic explanations that are largely absent in the smaller model.

\subsection{Discussion}

Overall, three conclusions emerge from the experiments.

First, scaling the model from picochat to nanochat substantially improves
language modeling quality and benchmark performance.

Second, the SwiGLU architecture slightly improves several reasoning
benchmarks but does not outperform the baseline in BPB and incurs
additional computational cost.

Third, architectural modifications can behave differently at different
scales, highlighting the importance of empirical evaluation when
designing transformer architectures.

\subsection{Codebase Organization}

This project consists of two main codebases: the modified \texttt{nanochat}
training framework used for model development and evaluation, and the
\texttt{AI-Accountant} repository that organizes the experimental workflow
for Assignment~3. The following section summarizes the key modifications and
directory structure used to support the experiments.

\subsubsection{Nanochat Modifications}

The nanochat framework was extended to support architectural ablations and
more robust evaluation. Compared to the original upstream repository, three
files were modified.

\begin{itemize}

\item \textbf{\texttt{nanochat/gpt.py}} \\
This file contains the core transformer architecture implementation.
Several extensions were introduced to support the architectural ablations
explored in this assignment. Specifically, the model now supports multiple
MLP variants (\texttt{relu2} and \texttt{swiglu}), optional residual
scaling through a LayerScale variant, and a configurable
\texttt{layerscale\_init} parameter controlling residual scaling strength.
Additionally, dynamic rotary embedding cache growth was implemented through
the \texttt{\_ensure\_rotary\_cache} function to allow evaluation with longer
sequence lengths than those used during training.

\item \textbf{\texttt{scripts/base\_train.py}} \\
The training entrypoint was modified to expose the new architectural
configuration options through the command-line interface and configuration
files. This allows experiments to select between the baseline ReLU$^2$ MLP,
SwiGLU MLP, and residual variants directly through experiment YAML files
without modifying the model implementation.

\item \textbf{\texttt{scripts/base\_eval.py}} \\
The evaluation pipeline was extended to support resumable CORE benchmark
evaluation. Progress files are written under a \texttt{base\_eval\_progress}
directory, enabling interrupted evaluation runs to resume without repeating
completed tasks. Additional features include task-level progress tracking,
support for \texttt{max-per-task} evaluation limits, and progress-path
reporting in logs for easier monitoring of long-running benchmark jobs.

\end{itemize}

These modifications enable flexible experimentation with architectural
variants while ensuring that long-running evaluation jobs can be resumed
safely.

\subsubsection{Assignment Repository Structure}


\definecolor{dircmt}{HTML}{6A737D}
\definecolor{filecmt}{HTML}{FF9966}
\newcommand{\cmt}[1]{\normalfont{\textcolor{dircmt}{\textit{#1}}}}
\newcommand{\fcmt}[1]{\normalfont{\textcolor{filecmt}{\textit{#1}}}}

How the \texttt{a3/} folder is organized inside the AI-Accountant repo is explained. The
structure follows the assignment's four-part layout.

\subsection*{Layout}

\begin{forest}
  for tree={
    font=\ttfamily,
    grow'=0,
    child anchor=west,
    parent anchor=south,
    anchor=west,
    calign=first,
    edge path={
      \noexpand\path [draw, -{|}] (!u.south west) +(7.5pt,0) |- node[fill,inner sep=1.25pt] {} (.child anchor)\forestoption{edge label};
    },
    before typesetting nodes={
      if n=1 {insert before={[,phantom]}} {}
    },
    fit=band,
    before computing xy={l=19.5pt},
    s sep=4pt,
  }
  [a3/
    [docs/ \cmt{teammate onboarding and shared references}]
    [shared/ \cmt{pipeline infrastructure used by all parts}
      [configs/ \cmt{smoke-test YAMLs for pipeline validation}]
      [scripts/ \cmt{Modal pipeline code}
        [main.py \fcmt{local entrypoint: reads YAML config, launches Modal run}]
        [shared/
          [infra.py \fcmt{Modal app, image, volume, and secrets definitions}]
          [helpers.py \fcmt{pure utilities: git checkout, checkpoint discovery, eval parsing}]
        ]
        [runners/
          [run.py \fcmt{orchestrator: dispatches training stages and evals in parallel}]
          [train.py \fcmt{runs nanochat training on a Modal GPU}]
          [evaluate.py \fcmt{runs BPB, CORE, and custom evals on a Modal GPU}]
        ]
      ]
    ]
    [p2/ \cmt{picochat ablation}
      [configs/
        [p2\_baseline.yaml \fcmt{vanilla picochat baseline}]
        [p2\_swiglu.yaml \fcmt{SwiGLU activation ablation}]
        [p2\_layerscale.yaml \fcmt{LayerScale ablation}]
        [p2\_swiglu\_layerscale.yaml \fcmt{combined SwiGLU + LayerScale ablation}]
        [p2\_*\_eval\_bpb\_core.yaml \fcmt{eval-only configs for each ablation}]
        [p2\_swiglu\_layerscale\_smoke.yaml \fcmt{quick smoke test for combined variant}]
      ]
      [training-troubleshooting.md \fcmt{notes on training issues encountered}]
    ]
    [p3/ \cmt{continued on next page}]
    [p4/ \cmt{continued on next page}]
  ]
\end{forest}

\newpage

\subsection*{Layout (continued)}

\begin{forest}
  for tree={
    font=\ttfamily,
    grow'=0,
    child anchor=west,
    parent anchor=south,
    anchor=west,
    calign=first,
    edge path={
      \noexpand\path [draw, -{|}] (!u.south west) +(7.5pt,0) |- node[fill,inner sep=1.25pt] {} (.child anchor)\forestoption{edge label};
    },
    before typesetting nodes={
      if n=1 {insert before={[,phantom]}} {}
    },
    fit=band,
    before computing xy={l=19.5pt},
    s sep=4pt,
  }
  [a3/
    [p3/ \cmt{context window extension}
      [docs/ \cmt{literature review, training config, eval design, final spec}]
      [configs/
        [p3\_baseline.yaml \fcmt{three-stage context extension experiment}]
        [p3\_eval\_only.yaml \fcmt{re-run evals on existing checkpoints}]
        [p3\_core\_evals.yaml \fcmt{CORE-only eval variant}]
        [p3\_from\_p2\_template.yaml \fcmt{template to eval P2's best model on P3 evals}]
      ]
      [evals/
        [context\_length\_eval.py \fcmt{positional perplexity eval on PG19 test split}]
      ]
      [results/ \cmt{run outputs}
        [core\_csvs/ \cmt{CORE per-task accuracy CSVs}]
        [custom\_evals/ \cmt{positional perplexity JSONs}]
        [figures/ \cmt{plots and summary tables}]
        [tex/ \cmt{compiled writeup PDF}]
      ]
      [analysis/
        [generate\_plots.py \fcmt{builds figures and tables from fetched results}]
      ]
    ]
    [p4/ \cmt{scaling laws and emergent abilities}
      [configs/
        [p4\_nano\_baseline\_fp8.yaml \fcmt{nanochat baseline with FP8 training}]
        [p4\_nano\_swiglu\_fp8.yaml \fcmt{nanochat SwiGLU variant with FP8 training}]
        [p4\_pipeline\_smoke\_20min.yaml \fcmt{quick smoke test for nano-scale runs}]
        [p4\_*\_eval\_only.yaml \fcmt{eval-only configs for each nano model}]
        [p4\_template.yaml \fcmt{template YAML for scaling runs}]
      ]
      [scripts/
        [compare\_emergent\_abilities.py \fcmt{finds questions nanochat answers but picochat cannot}]
      ]
      [results/ \cmt{run outputs}
        [final/ \cmt{final eval CSVs and JSONs per model}]
      ]
      [emergent\_abilities\_results.txt \fcmt{emergent abilities comparison output}]
    ]
  ]
\end{forest}


\newpage
\paragraph{Part 1: Infrastructure and Orchestration}

The infrastructure and orchestration layer is implemented under
\texttt{a3/shared}. This directory contains reusable pipeline components
used by all experiments.

Key files include:

\begin{itemize}

\item \textbf{\texttt{scripts/main.py}} \\
Primary local entrypoint for running experiments from YAML configuration
files.

\item \textbf{\texttt{scripts/runners/run.py}} \\
Orchestrates the full training and evaluation pipeline by coordinating
individual stages.

\item \textbf{\texttt{scripts/runners/train.py}} \\
Implements the training worker responsible for launching model training
jobs.

\item \textbf{\texttt{scripts/runners/evaluate.py}} \\
Handles evaluation runs, including CORE benchmark execution.

\item \textbf{\texttt{scripts/shared/infra.py}} \\
Defines the Modal infrastructure configuration, including container
images, persistent volumes, and secret management.

\item \textbf{\texttt{scripts/shared/helpers.py}} \\
Provides helper utilities for parsing configuration files and managing
pipeline stages.

\end{itemize}

The \texttt{a3/docs} directory contains supporting documentation for the
infrastructure design, including interface definitions and experiment
workflow descriptions.

\paragraph{Part 2: Picochat Ablation Experiments}

All pico-scale ablation experiments are defined in
\texttt{a3/p2/configs}. These configuration files specify experiment
variants including:

\begin{itemize}
\item baseline architecture
\item SwiGLU MLP variant
\item LayerScale residual variant
\item combined architectural modifications
\item smoke-test and evaluation-only configurations
\end{itemize}

The file \texttt{training-troubleshooting.md} documents issues encountered
during early training runs and the steps taken to resume interrupted
experiments.

\paragraph{Part 3: Analysis Structure}

The \texttt{a3/p3} directory provides the structure for analysis artifacts,
including placeholders for evaluation results, documentation, and
experiment outputs. While this directory contains minimal code, it
establishes the standardized layout used for subsequent parts of the
assignment.

\paragraph{Part 4: Nanochat Experiments}

The final nanochat experiments are organized under \texttt{a3/p4}.

\begin{itemize}

\item \textbf{\texttt{configs}} \\
Contains YAML configurations for full nanochat training runs, evaluation-only
jobs, hyperparameter sweeps, and smoke-test pipelines.

\item \textbf{\texttt{results}} \\
Stores downloaded evaluation artifacts including BPB measurements,
CORE benchmark outputs, and aggregated CSV/JSON summaries used in the
analysis.

\item \textbf{\texttt{scripts/compare\_emergent\_abilities.py}} \\
A diagnostic script that generates side-by-side responses from picochat
and nanochat models to compare qualitative reasoning abilities.

\item \textbf{\texttt{emergent\_abilities\_results.txt}} \\
A generated file containing the qualitative comparison outputs used in
the emergent ability analysis in Part~4.

\end{itemize}

Overall, this structure separates reusable infrastructure code, experiment
configuration files, and analysis artifacts. This organization improves
reproducibility and makes it easier to track experimental results across
different parts of the assignment.